\DocumentMetadata{
lang = en-US,
pdfversion = 2.0,
pdfstandard = ua-2,
tagging = on,
tagging-setup = {math/setup=mathml-SE}
}

\documentclass[12pt,oneside]{book}

%accessibility stuff
\usepackage{tagpdf}
\tagpdfsetup{activate-all}
\tagpdfsetup{math/alt/use}

\usepackage{caption}
\usepackage[labelfont={}]{subcaption}
\captionsetup{subrefformat=parens}
\usepackage{epigraph}
\usepackage[table]{xcolor}
% IBM palette verified at https://davidmathlogic.com/colorblind/
\definecolor{IBMblue}{HTML}{648FFF}
\definecolor{IBMpurple}{HTML}{785EF0}
\definecolor{IBMmagenta}{HTML}{DC267F}
\definecolor{IBMorange}{HTML}{FE6100}
\definecolor{IBMgold}{HTML}{FFB000}

\colorlet{blue}{IBMblue}
\colorlet{pink}{IBMmagenta}
\colorlet{magenta}{IBMpurple}
\colorlet{green}{IBMblue}
\colorlet{red}{IBMmagenta}
\colorlet{cyan}{IBMblue}
\colorlet{yellow}{IBMgold}
\usepackage{tikz}
\usetikzlibrary{arrows,automata,positioning,fit,shapes,matrix}
\usetikzlibrary{decorations.markings}
\usetikzlibrary{decorations.pathreplacing}

% Shared edge thickness for the selected figures.
\tikzset{figure edge/.style={line width=0.8pt}}
\tikzset{colored fill/.style={fill opacity=0.85,text opacity=1}}
% Target border width at the common 3 cm display height.
\newlength{\solidborderwidth}
\setlength{\solidborderwidth}{0.6pt}
% Compensate for each drawing's native size before resizebox.
\tikzset{
rhombic border/.style={draw=black,line width=1.386\solidborderwidth},
heap border/.style={draw=black,line width=1.110\solidborderwidth},
tromino border/.style={draw=black,line width=2.017\solidborderwidth}
}

% Shared vertex styles.
\tikzset{
graph vertex/.style={draw,circle,inner sep=1.25pt},
quotient vertex/.style={circle,draw,fill=grey,inner sep=0pt,minimum size=3mm}
}

% Shared styles for gamegraph figures.
\tikzset{
rect vert/.style={rectangle,rounded corners,draw,fill=grey!40,inner sep=2.7pt,minimum size=4mm},
d/.style={draw,thick,black,-stealth}
}

\usepackage{kpfonts-otf}
\usepackage{amsmath}
% \usepackage{chessfss}
\usepackage{amsthm}
\usepackage[letterpaper,landscape,margin=.65in]{geometry}
\usepackage[hang,flushmargin]{footmisc}
\usepackage{color}
\colorlet{darkblue}{IBMblue}
\definecolor{grey}{rgb}{.7, .7, .7}
\colorlet{darkblue}{IBMblue}
\definecolor{grey}{rgb}{.7, .7, .7}
\colorlet{rred}{IBMmagenta}
\colorlet{darkred}{IBMmagenta!80!black}
\colorlet{rose}{IBMmagenta}
\colorlet{ggreen}{IBMblue}
\colorlet{naugreen}{IBMblue}
\colorlet{springgreen}{IBMblue!70}
\colorlet{darkblue}{IBMblue}
\colorlet{turq}{IBMgold}
\colorlet{naublue}{IBMblue}
\colorlet{mediterranean}{IBMpurple}
\colorlet{butterfly}{IBMblue}
\colorlet{icyblue}{IBMblue}
\colorlet{orange2}{IBMorange}
\colorlet{orange}{IBMgold}
\colorlet{nectarine}{IBMorange}
\colorlet{purple}{IBMpurple}
\colorlet{darkorchid}{IBMorange}
\colorlet{purple2}{IBMpurple}
\colorlet{NAUTrueBlue}{black}
\colorlet{NAUGold}{IBMgold}
\colorlet{Supai}{IBMblue}
\colorlet{RedRock}{IBMmagenta}
\colorlet{Monsoon}{IBMgold}
\colorlet{SummerShade}{IBMpurple}
\colorlet{Sky}{IBMblue}
\colorlet{Sunset}{IBMorange}

\usepackage[most]{tcolorbox}
\usepackage[breaklinks]{hyperref}
\hypersetup{
pdfsubject={An inquiry-based problem sequence for a graduate course in enumerative combinatorics},
pdfkeywords={enumerative combinatorics, counting, bijections, generating functions, combinatorial species, inquiry-based learning}
}
\hypersetup{
colorlinks=true,
linkcolor=darkblue,
citecolor=darkblue,
urlcolor=darkblue,
bookmarksnumbered
}

% Restore annotation underlines after colorlinks enables colored text.
\hypersetup{
pdfborderstyle={/S/U/W 0.5},
linkbordercolor=darkblue,
citebordercolor=darkblue,
urlbordercolor=darkblue,
filebordercolor=darkblue
}

%\setlength{\parindent}{0pt}

\newsavebox{\savepar}
\newenvironment{textbox}{\noindent\begin{lrbox}{\savepar}\begin{minipage}[c]{.98\textwidth}}{\end{minipage}\end{lrbox}\fcolorbox{black}{white}{\usebox{\savepar}}}

\usetikzlibrary{3d}
\usetikzlibrary{arrows,shapes,positioning,backgrounds}
\usetikzlibrary{decorations.markings}
\usetikzlibrary{decorations.pathreplacing}
\usetikzlibrary{shapes.geometric}
\usetikzlibrary{calc}
\pgfdeclarelayer{background}
\pgfdeclarelayer{gamebackground}
\pgfsetlayers{background,gamebackground,main}

\tikzset{
gamepos/.style={rectangle,rounded corners,draw,fill=grey!40,
inner sep=3pt,minimum height=6mm},
gamearrow/.style={draw,thick,black,-stealth},
slidepos/.style={gamepos,inner xsep=2pt,inner ysep=1.5pt,
minimum height=5mm,font=\small},
trebpos/.style={slidepos},
trebarrow/.style={draw,semithick,black,-stealth},
qzero/.style={fill=IBMorange!65},
qone/.style={fill=black!18},
qtwo/.style={fill=IBMblue!45},
qthree/.style={fill=IBMgold!65},
qfour/.style={fill=IBMpurple!40},
qfive/.style={fill=IBMgold!25},
qsix/.style={fill=IBMmagenta!45}
}

%for highlighting paths
\tikzset{
my box/.style = {
, line cap = round
, line join = round
}
}

\newcommand{\highlight}[3]{
\path [my box, line width = #1, draw = #2,opacity=.2] #3;
}

\newcommand{\convexpath}[2]{
[
create hullnodes/.code={
\global\edef\namelist{#1}
\foreach [count=\counter] \nodename in \namelist {
\global\edef\numberofnodes{\counter}
\node at (\nodename) [draw=none,name=hullnode\counter] {};
}
\node at (hullnode\numberofnodes) [name=hullnode0,draw=none] {};
\pgfmathtruncatemacro\lastnumber{\numberofnodes+1}
\node at (hullnode1) [name=hullnode\lastnumber,draw=none] {};
},
create hullnodes
]
($(hullnode1)!#2!-90:(hullnode0)$)
\foreach [
evaluate=\currentnode as \previousnode using \currentnode-1,
evaluate=\currentnode as \nextnode using \currentnode+1
] \currentnode in {1,...,\numberofnodes} {
let
\p1 = ($(hullnode\currentnode)!#2!-90:(hullnode\previousnode)$),
\p2 = ($(hullnode\currentnode)!#2!90:(hullnode\nextnode)$),
\p3 = ($(\p1) - (hullnode\currentnode)$),
\n1 = {atan2(\y3,\x3)},
\p4 = ($(\p2) - (hullnode\currentnode)$),
\n2 = {atan2(\y4,\x4)},
\n{delta} = {-Mod(\n1-\n2,360)}
in
{-- (\p1) arc[start angle=\n1, delta angle=\n{delta}, radius=#2] -- (\p2)}
}
-- cycle
}

\setlength{\epigraphwidth}{0.5\textwidth}
\renewcommand{\sourceflush}{flushleft}

\setlength{\parindent}{0pt}

\newcommand{\blankline}{\pagebreak[2]\vspace{.5\baselineskip}}

\tcbset{size=fbox}


\usepackage{graphicx}
\usepackage{longtable}
\newlength{\figureheight}
\setlength{\figureheight}{3cm}
\newsavebox{\figurecellbox}
\newcommand{\figurecell}[2]{%
\begin{minipage}[t]{.47\textwidth}
\centering
\sbox{\figurecellbox}{#2}%
\resizebox*{!}{\figureheight}{\usebox{\figurecellbox}}\par
\vspace{5pt}
{\small #1\par}
\end{minipage}%
}

\begin{document}
\begin{center}
{\Large Figure comparison}\par
\small All diagrams are 3 cm high; descriptions are excluded from that height.
\end{center}
\setlength{\tabcolsep}{0pt}
\begin{longtable}{@{}c@{\hspace{.04\textwidth}}c@{}}
\figurecell{Fraenkel tree}{%
\begin{tikzpicture}[x=1.8em,y=4ex,scale=0.8]
\begin{scope}[every node/.style={graph vertex}]
\node (a) at (0,0) {};
\node (b) at (1,0) {};
\node (c) at (2,0) {};
\node (d) at (3,0) {};
\node (e) at (4,0) {};
\node (v1_0) at (0,1) {};
\node (v1_1) at (1.5,1) {};
\node (v1_2) at (3,1) {};
\node (v1_3) at (4,1) {};
\node (v2_0) at (0,2) {};
\node (v2_1) at (1.5,2) {};
\node (v2_2) at (3.5,2) {};
\node (v3_0) at (0,3) {};
\node (v3_1) at (2.5,3) {};
\node (v4_0) at (2,4) {};
\end{scope}
\draw[figure edge](v1_0) -- (a);
\draw[figure edge](v1_1) -- (b);
\draw[figure edge](v1_1) -- (c);
\draw[figure edge](v1_2) -- (d);
\draw[figure edge](v1_3) -- (e);
\draw[figure edge](v2_0) -- (v1_0);
\draw[figure edge](v2_1) -- (v1_1);
\draw[figure edge](v2_2) -- (v1_2);
\draw[figure edge](v2_2) -- (v1_3);
\draw[figure edge](v3_0) -- (v2_0);
\draw[figure edge](v3_1) -- (v2_1);
\draw[figure edge](v3_1) -- (v2_2);
\draw[figure edge](v4_0) -- (v3_0);
\draw[figure edge](v4_0) -- (v3_1);
\end{tikzpicture}
} &
\figurecell{Rhombic tilings}{%
{
\newcommand{\pptop}[3]{
\path[colored fill,fill=IBMgold,rhombic border]
(#1,#2,#3) -- ({#1+1},#2,#3) -- ({#1+1},{#2+1},#3) -- (#1,{#2+1},#3) -- cycle;
}
\newcommand{\ppred}[3]{
\path[colored fill,fill=IBMblue,rhombic border]
(#1,#2,#3) -- (#1,{#2+1},#3) -- (#1,{#2+1},{#3+1}) -- (#1,#2,{#3+1}) -- cycle;
}
\newcommand{\ppblue}[3]{
\path[colored fill,fill=IBMmagenta,rhombic border]
(#1,#2,#3) -- ({#1+1},#2,#3) -- ({#1+1},#2,{#3+1}) -- (#1,#2,{#3+1}) -- cycle;
}
\newcommand{\pptiletop}[3]{
\path[fill=white,rhombic border]
(#1,#2,#3) -- ({#1+1},#2,#3) -- ({#1+1},{#2+1},#3) -- (#1,{#2+1},#3) -- cycle;
}
\newcommand{\pptilex}[3]{
\path[fill=white,rhombic border]
(#1,#2,#3) -- (#1,{#2+1},#3) -- (#1,{#2+1},{#3+1}) -- (#1,#2,{#3+1}) -- cycle;
}
\newcommand{\pptiley}[3]{
\path[fill=white,rhombic border]
(#1,#2,#3) -- ({#1+1},#2,#3) -- ({#1+1},#2,{#3+1}) -- (#1,#2,{#3+1}) -- cycle;
}
\begin{tikzpicture}[line join=round,line cap=round]
\begin{scope}[shift={(-3.4,0)}]
\begin{scope}[x={(0.68cm,-0.34cm)},y={(-0.68cm,-0.34cm)},z={(0cm,0.68cm)}]
% Complete the hexagon with seven white rhombi outlined in black.
\begin{scope}
\colorlet{NAUTrueBlue}{black}
\pptilex{0}{1}{2}
\pptilex{0}{2}{1}
\pptilex{0}{2}{2}
\pptiley{2}{0}{2}
\pptiletop{2}{1}{0}
\pptiletop{1}{2}{0}
\pptiletop{2}{2}{0}
\end{scope}
% Red faces: positive x direction.
\ppred{2}{0}{2}
\ppred{3}{0}{0}
\ppred{3}{0}{1}
\ppred{1}{1}{1}
\ppred{2}{1}{0}
\ppred{1}{2}{0}
% Blue faces: positive y direction.
\ppblue{0}{1}{2}
\ppblue{0}{2}{1}
\ppblue{0}{3}{0}
\ppblue{1}{1}{1}
\ppblue{1}{1}{2}
\ppblue{1}{2}{0}
\ppblue{2}{1}{0}
\ppblue{2}{1}{1}
% Top faces of the six nonzero columns.
\pptop{0}{0}{3}
\pptop{1}{0}{3}
\pptop{2}{0}{2}
\pptop{0}{1}{2}
\pptop{1}{1}{1}
\pptop{0}{2}{1}
% Outer hexagon.
\draw[rhombic border]
(0,0,3) -- (3,0,3) -- (3,0,0) -- (3,3,0) -- (0,3,0) -- (0,3,3) -- cycle;
\end{scope}
\end{scope}
\end{tikzpicture}
}
} \\[12pt]

\figurecell{Graph of a permutation}{%
\begin{tikzpicture}[alt={A seven by seven grid showing the 231-avoiding permutation 2173546. Columns run from left to right and rows from bottom to top. A dot is centered in each column's square at row 2, 1, 7, 3, 5, 4, and 6, respectively. The maximum occurs in column 3.},x=0.65cm,y=0.65cm]
\draw[step=1,gray!25,very thin] (0,0) grid (7,7);
\foreach \i/\j in {1/2,2/1,3/7,4/3,5/5,6/4,7/6}
\fill ({\i-0.5},{\j-0.5}) circle (4pt);
\end{tikzpicture}
} &
\figurecell{Dyck path}{%
% Bottom-hugging staircase: column i has height min(w(i),...,w(7))-1.
% For 2173546 these heights are 0,0,2,2,3,3,5; reflect across y=x.
\begin{tikzpicture}[alt={The Dyck path corresponding to the permutation 2173546, obtained by reflecting its bottom-hugging staircase across y equals x. On a seven by seven grid it follows north, north, east, east, north, north, east, north, north, east, east, north, east, east, staying on or above the diagonal.},x=0.65cm,y=0.65cm]
\draw[step=1,gray!25,very thin] (0,0) grid (7,7);
\draw[black,line width=1.8pt,line cap=round,line join=round]
(0,0) -- (0,2) -- (2,2) -- (2,4) -- (3,4)
-- (3,6) -- (5,6) -- (5,7) -- (7,7);
\end{tikzpicture}
} \\[12pt]

\figurecell{Braid graph}{%
\begin{tikzpicture}[every circle node/.style={graph vertex}, scale = 1.2]
\node [circle] (1) at (0,1){};
\node [circle] (2) at (.707,.293){};
\node [circle] (3) at (-.707,1.707){};
\node [circle] (4) at (.707,1.707){};
\node [circle] (5) at (0,2.414){};
\draw[rotate = 90, pink,-, very thick,figure edge,line width=1.4pt] (1) to (2);
\draw[pink,-, very thick,figure edge,line width=1.4pt] (1) to (3);
\draw[black,-, very thick,figure edge] (1) to (4);
\draw[black,-, very thick,figure edge] (3) to (5);
\draw[blue,-, very thick,figure edge,line width=1.4pt] (5) to (4);
%%Flapperbox
\node [circle] (9) at (1.414,1){};
\node [circle] (10) at (0,1.707){};
\node [circle] (11) at (0,3.121){};
\node [circle] (13) at (.707,2.414){};
\node [circle] (14) at (-.707,2.414){};
\node [circle] (15) at (1.414,3.121){};
\node [circle] (16) at (.707,3.828){};
\draw[pink,-, very thick,figure edge,line width=1.4pt] (1) to (10);
\draw[black,-, very thick,figure edge] (10) to (13);
\draw[pink,-, very thick,figure edge,line width=1.4pt] (10) to (14);
\draw[black,-, very thick,figure edge] (14) to (11);
\draw[blue,-, very thick,figure edge,line width=1.4pt] (13) to (11);
\draw[blue,-, very thick,figure edge,line width=1.4pt] (13) to (4);
\draw[blue,-, very thick,figure edge,line width=1.4pt] (4) to (9);
\draw[black,-, very thick,figure edge] (9) to (2);
\draw[blue,-, very thick,figure edge,line width=1.4pt] (5) to (11);
\draw[pink,-, very thick,figure edge,line width=1.4pt] (14) to (3);
\draw[blue,-, very thick,figure edge,line width=1.4pt] (15) to (16);
\draw[blue,-, very thick,figure edge,line width=1.4pt] (11) to (16);
\draw[blue,-, very thick,figure edge,line width=1.4pt] (15) to (13);
\begin{pgfonlayer}{background}
\highlight{8pt}{blue}{(9.center) to (4.center) to (5.center) to (11.center) to (16.center) to (15.center) to (13.center) to (11.center)}
\highlight{8pt}{blue}{(13.center) to (4.center)}
\highlight{8pt}{pink}{(2.center) to (1.center) to (3.center) to (14.center) to (10.center) to (1.center)}
\fill[gray,opacity=0.2] \convexpath{9, 4, 5, 11, 13,4}{8pt};
\end{pgfonlayer}
\end{tikzpicture}
} &
\figurecell{Permutahedron}{%
\begin{tikzpicture}[every circle node/.style={graph vertex}, scale=1]
% Nodes
\node [circle] (a) at (9,0) {};
\node [circle] (b) at (10,0) {};
\node [circle] (c) at (9.25,0.25) {};
\node [circle] (d) at (10.25,0.25) {};
\node [circle] (e) at (8,0.75) {};
\node [circle] (f) at (10.25,0.75) {};
\node [circle] (g) at (9,1.25) {};
\node [circle] (h) at (11,1) {};
\node [circle] (i) at (7.75,1.75) {};
\node [circle] (j) at (8.25,1.5) {};
\node [circle] (k) at (9.5,1.5) {};
\node [circle] (l) at (11,1.5) {};
\node [circle] (m) at (11.25,1.75) {};
\node [circle] (n) at (8,2.5) {};
\node [circle] (o) at (8.25,2) {};
\node [circle] (p) at (9,2.75) {};
\node [circle] (q) at (9.75,2) {};
\node [circle] (r) at (10.75,2) {};
\node [circle] (s) at (11,2.75) {};
\node [circle] (t) at (10.25,2.25) {};
\node [circle] (u) at (8.75,3.25) {};
\node [circle] (v) at (9,3.5) {};
\node [circle] (w) at (9.75,3.25) {};
\node [circle] (x) at (10,3.5) {};
% Orange Edges
\draw[IBMpurple,-, very thick,figure edge] (a) to (b);
\draw[IBMpurple,-, very thick,figure edge] (b) to (d);
\draw[IBMpurple,-, very thick, opacity=0.5,figure edge] (a) to (c);
\draw[IBMpurple,-, very thick, opacity=0.5,figure edge] (c) to (d);
\draw[IBMpurple,-, very thick,figure edge] (a) to (e);
\draw[IBMpurple,-, very thick,figure edge] (b) to (f);
\draw[IBMpurple,-, very thick,figure edge] (d) to (h);
\draw[IBMpurple,-, very thick, opacity=0.5,figure edge] (c) to (g);
\draw[IBMpurple,-, very thick,figure edge] (e) to (i);
\draw[IBMpurple,-, very thick,figure edge] (e) to (j);
\draw[IBMpurple,-, very thick,figure edge] (k) to (f);
\draw[IBMpurple,-, very thick,figure edge] (l) to (f);
\draw[IBMpurple,-, very thick,figure edge] (l) to (m);
\draw[IBMpurple,-, very thick,figure edge] (h) to (m);
\draw[IBMpurple,-, very thick,figure edge] (n) to (i);
\draw[IBMpurple,-, very thick,figure edge] (n) to (j);
\draw[IBMpurple,-, very thick, opacity=0.5,figure edge] (o) to (i);
\draw[IBMpurple,-, very thick, opacity=0.5,figure edge] (g) to (o);
\draw[IBMpurple,-, very thick, opacity=0.5,figure edge] (o) to (p);
\draw[IBMpurple,-, very thick, opacity=0.5,figure edge] (q) to (p);
\draw[IBMpurple,-, very thick, opacity=0.5,figure edge] (q) to (g);
\draw[IBMpurple,-, very thick, opacity=0.5,figure edge] (q) to (r);
\draw[IBMpurple,-, very thick, opacity=0.5,figure edge] (r) to (s);
\draw[IBMpurple,-, very thick, opacity=0.5,figure edge] (r) to (h);
\draw[IBMpurple,-, very thick,figure edge] (j) to (k);
\draw[IBMpurple,-, very thick, opacity=0.5,figure edge] (p) to (v);
\draw[IBMpurple,-, very thick,figure edge] (u) to (v);
\draw[IBMpurple,-, very thick,figure edge] (u) to (w);
\draw[IBMpurple,-, very thick,figure edge] (x) to (w);
\draw[IBMpurple,-, very thick,figure edge] (u) to (n);
\draw[IBMpurple,-, very thick,figure edge] (w) to (t);
\draw[IBMpurple,-, very thick,figure edge] (s) to (x);
\draw[IBMpurple,-, very thick,figure edge] (v) to (x);
\draw[IBMpurple,-, very thick,figure edge] (s) to (m);
\draw[IBMpurple,-, very thick,figure edge] (t) to (k);
\draw[IBMpurple,-, very thick,figure edge] (t) to (l);
\end{tikzpicture}
} \\[12pt]

\figurecell{Box product of graphs}{%
\begin{tikzpicture}[every circle node/.style={graph vertex}]
\node [circle] (1) at (1.3,0){};
\node [circle] (2) at ({1.3*cos(45)},{1.3*sin(45)}){};
\node [circle] (3) at (0,1.3){};
\node [circle] (4) at ({1.3*cos(135)},{1.3*sin(135)}){};
\node [circle] (5) at (-1.3,0){};
\node [circle] (6) at ({1.3*cos(225)},{1.3*sin(225)}){};
\node [circle] (7) at (0,-1.3){};
\node [circle] (8) at ({1.3*cos(315)},{1.3*sin(315)}){};
\node [circle] (9) at (.55,0){};
\node [circle] (10) at ({.55*cos(45)},{.55*sin(45)}){};
\node [circle] (11) at (0,.55){};
\node [circle] (12) at ({.55*cos(135)},{.55*sin(135)}){};
\node [circle] (13) at (-.55,0){};
\node [circle] (14) at ({.55*cos(225)},{.55*sin(225)}){};
\node [circle] (15) at (0,-.55){};
\node [circle] (16) at ({.55*cos(315)},{.55*sin(315)}){};
\node [circle] (17) at (1.3+.92,-.393){};
\node [circle] (18) at (0+.92,-1.3-.393){};
\node [circle] (19) at ({1.3*cos(315)+.92},{1.3*sin(315)-.393}){};
\node [circle] (20) at ({.55*cos(315)+.92},{.55*sin(315)-.393}){};
\node [circle] (21) at (-1.3-.92,0-.393){};
\node [circle] (22) at ({1.3*cos(225)-.92},{1.3*sin(225)-.393}){};
\node [circle] (23) at (0-.92,-1.3-.393){};
\node [circle] (24) at ({.55*cos(225)-.92},{.55*sin(225)-.393}){};
\node [circle] (25) at (0,-2.05){};
\draw[magenta,-, very thick,figure edge] (1) to (2);
\draw[magenta,-, very thick,figure edge] (2) to (3);
\draw[turq,-, very thick,figure edge] (3) to (4);
\draw[turq,-, very thick,figure edge] (4) to (5);
\draw[magenta,-, very thick,figure edge] (5) to (6);
\draw[magenta,-, very thick,figure edge] (6) to (7);
\draw[turq,-, very thick,figure edge] (7) to (8);
\draw[turq,-, very thick,figure edge] (8) to (1);
\draw[magenta,-, very thick,figure edge] (8) to (9);
\draw[turq,-, very thick,figure edge] (1) to (16);
\draw[turq,-, very thick,figure edge] (2) to (9);
\draw[magenta,-, very thick,figure edge] (3) to (10);
\draw[turq,-, very thick,figure edge] (5) to (12);
\draw[turq,-, very thick,figure edge] (6) to (13);
\draw[magenta,-, very thick,figure edge] (8) to (15);
\draw[magenta,-, very thick,figure edge] (1) to (10);
\draw[turq,-, very thick,figure edge] (2) to (11);
\draw[turq,-, very thick,figure edge] (3) to (12);
\draw[magenta,-, very thick,figure edge] (4) to (13);
\draw[magenta,-, very thick,figure edge] (5) to (14);
\draw[turq,-, very thick,figure edge] (6) to (15);
\draw[turq,-, very thick,figure edge] (7) to (16);
\draw[magenta,-, very thick,figure edge] (9) to (12);
\draw[turq,-, very thick,figure edge] (10) to (13);
\draw[magenta,-, very thick,figure edge] (12) to (15);
\draw[magenta,-, very thick,figure edge] (13) to (16);
\draw[turq,-, very thick,figure edge] (14) to (9);
\draw[turq,-, very thick,figure edge] (15) to (10);
\draw[magenta,-, very thick,figure edge] (16) to (11);
\draw[turq,-, very thick,figure edge] (22) to (6);
\draw[turq,-, very thick,figure edge] (21) to (5);
\draw[turq,-, very thick,figure edge] (23) to (7);
\draw[turq,-, very thick,figure edge] (24) to (14);
\draw[magenta,-, very thick,figure edge] (19) to (8);
\draw[magenta,-, very thick,figure edge] (17) to (1);
\draw[magenta,-, very thick,figure edge] (18) to (7);
\draw[magenta,-, very thick,figure edge] (20) to (16);
\draw[turq,-, very thick,figure edge] (17) to (19);
\draw[turq,-, very thick,figure edge] (18) to (19);
\draw[turq,-, very thick,figure edge] (17) to (20);
\draw[turq,-, very thick,figure edge] (20) to (18);
\draw[magenta,-, very thick,figure edge] (21) to (22);
\draw[magenta,-, very thick,figure edge] (22) to (23);
\draw[magenta,-, very thick,figure edge] (23) to (24);
\draw[magenta,-, very thick,figure edge] (24) to (21);
\draw[magenta,-, very thick,figure edge] (7) to (14);
\draw[turq,-, very thick,figure edge] (11) to (14);
\draw[magenta,-, very thick,figure edge] (4) to (11);
\draw[turq,-, very thick,figure edge] (18) to (25);
\draw[magenta,-, very thick,figure edge] (23) to (25);
\end{tikzpicture}
} &
\figurecell{Hyperplane arrangement}{%
\begin{tikzpicture}[scale=0.88,transform shape]
\begin{scope}[shift={(-5.25,-1.85)},x={(1cm,0cm)},y={(0.32cm,0.22cm)},z={(0cm,1cm)}]
\filldraw[fill=purple2,fill opacity=0.09,draw=purple2!70!black,draw opacity=0.45,line width=0.45pt]
(1.6,0,0) -- (1.6,3.2,0) -- (1.6,3.2,3.2) -- (1.6,0,3.2) -- cycle;
\filldraw[fill=purple2,fill opacity=0.09,draw=purple2!70!black,draw opacity=0.45,line width=0.45pt]
(0,1.6,0) -- (3.2,1.6,0) -- (3.2,1.6,3.2) -- (0,1.6,3.2) -- cycle;
\filldraw[fill=purple2,fill opacity=0.09,draw=purple2!70!black,draw opacity=0.45,line width=0.45pt]
(0,0,1.6) -- (3.2,0,1.6) -- (3.2,3.2,1.6) -- (0,3.2,1.6) -- cycle;
\filldraw[fill=purple2,fill opacity=0.09,draw=purple2!70!black,draw opacity=0.45,line width=0.45pt]
(0,0,0) -- (3.2,3.2,0) -- (3.2,3.2,3.2) -- (0,0,3.2) -- cycle;
\filldraw[fill=purple2,fill opacity=0.09,draw=purple2!70!black,draw opacity=0.45,line width=0.45pt]
(0,3.2,0) -- (3.2,0,0) -- (3.2,0,3.2) -- (0,3.2,3.2) -- cycle;
\filldraw[fill=purple2,fill opacity=0.09,draw=purple2!70!black,draw opacity=0.45,line width=0.45pt]
(0,0,0) -- (3.2,0,3.2) -- (3.2,3.2,3.2) -- (0,3.2,0) -- cycle;
\filldraw[fill=purple2,fill opacity=0.09,draw=purple2!70!black,draw opacity=0.45,line width=0.45pt]
(0,0,3.2) -- (3.2,0,0) -- (3.2,3.2,0) -- (0,3.2,3.2) -- cycle;
\filldraw[fill=purple2,fill opacity=0.09,draw=purple2!70!black,draw opacity=0.45,line width=0.45pt]
(0,0,0) -- (3.2,0,0) -- (3.2,3.2,3.2) -- (0,3.2,3.2) -- cycle;
\filldraw[fill=purple2,fill opacity=0.09,draw=purple2!70!black,draw opacity=0.45,line width=0.45pt]
(0,0,3.2) -- (3.2,0,3.2) -- (3.2,3.2,0) -- (0,3.2,0) -- cycle;
\draw[black,line width=1.2pt,line cap=round,line join=round]
(0,0,0) -- (3.2,0,0) -- (3.2,3.2,0) -- (0,3.2,0) -- cycle
(0,0,3.2) -- (3.2,0,3.2) -- (3.2,3.2,3.2) -- (0,3.2,3.2) -- cycle
(0,0,0) -- (0,0,3.2)
(3.2,0,0) -- (3.2,0,3.2)
(3.2,3.2,0) -- (3.2,3.2,3.2)
(0,3.2,0) -- (0,3.2,3.2);
\end{scope}
\end{tikzpicture}
} \\[12pt]

\figurecell{Game graph}{%
\begin{tikzpicture}[x=.86cm,y=.90cm]
\coordinate (t0000) at (0,5.35);
\coordinate (t1000) at (-2.7,4.05);
\coordinate (t0100) at (-.9,4.05);
\coordinate (t0010) at (.9,4.05);
\coordinate (t0001) at (2.7,4.05);
\coordinate (t1100) at (-3.4,2.65);
\coordinate (t1010) at (-2.04,2.65);
\coordinate (t1001) at (-.68,2.65);
\coordinate (t0110) at (.68,2.65);
\coordinate (t0101) at (2.04,2.65);
\coordinate (t0011) at (3.4,2.65);
\coordinate (t1110) at (-2.35,1.25);
\coordinate (t1101) at (-.78,1.25);
\coordinate (t1011) at (.78,1.25);
\coordinate (t0111) at (2.35,1.25);
\coordinate (t1111) at (0,-.15);
\node[trebpos,qsix] (p0000) at (t0000) {$0000$};
\node[trebpos,qfive] (p1000) at (t1000) {$1000$};
\node[trebpos,qfour] (p0100) at (t0100) {$0100$};
\node[trebpos,qfour] (p0010) at (t0010) {$0010$};
\node[trebpos,qfive] (p0001) at (t0001) {$0001$};
\node[trebpos,qthree] (p1100) at (t1100) {$1100$};
\node[trebpos,qthree] (p1010) at (t1010) {$1010$};
\node[trebpos,qtwo] (p1001) at (t1001) {$1001$};
\node[trebpos,qone] (p0110) at (t0110) {$0110$};
\node[trebpos,qthree] (p0101) at (t0101) {$0101$};
\node[trebpos,qthree] (p0011) at (t0011) {$0011$};
\node[trebpos,qzero] (p1110) at (t1110) {$1110$};
\node[trebpos,qone] (p1101) at (t1101) {$1101$};
\node[trebpos,qone] (p1011) at (t1011) {$1011$};
\node[trebpos,qzero] (p0111) at (t0111) {$0111$};
\node[trebpos,qzero] (p1111) at (t1111) {$1111$};
% Named nodes clip each arrow at the cell boundary; the background layer
% keeps all connectors behind the filled cells.
\begin{pgfonlayer}{gamebackground}
\path[trebarrow,figure edge]
(p0000) edge (p1000)
edge (p0100)
edge (p0010)
edge (p0001)
(p1000) edge (p1100)
edge (p1010)
edge (p1001)
(p0100) edge (p1100)
edge (p0110)
edge (p0101)
(p0010) edge (p1010)
edge (p0110)
edge (p0011)
(p0001) edge (p1001)
edge (p0101)
edge (p0011)
(p1100) edge (p1110)
edge (p1101)
(p1010) edge (p1110)
edge (p1011)
(p1001) edge (p1101)
edge (p1011)
(p0110) edge (p1110)
edge (p0111)
(p0101) edge (p1101)
edge (p0111)
(p0011) edge (p1011)
edge (p0111)
(p1101) edge (p1111)
(p1011) edge (p1111);
\end{pgfonlayer}
\end{tikzpicture}
} &
\figurecell{Type II heap}{%
% Copied from chapter1.tex, Figure 1.13, lower-right heap.
% Helpers are local to this figure.
{%
\begin{tikzpicture}
[x={(-0.5cm,-0.5cm)}, y={(1cm,0cm)}, z={(0cm,1cm)},
scale=0.3,every node/.append style={thick, transform shape, scale=2}]
\newcommand\drawface{\draw[heap border,colored fill,fill=heapcolor] (-1,-1) rectangle (1,1)}
\newcommand\faceone{
\begin{scope}[canvas is yz plane at x=-1]
\drawface;
\end{scope}
}
\newcommand\facetwo{
\begin{scope}[canvas is yx plane at z=-1]
\drawface;
\end{scope}
}
\newcommand\facethree{
\begin{scope}[canvas is zx plane at y=-1]
\drawface;
\end{scope}
}
\newcommand\rightonebar{
\begin{scope}[canvas is zx plane at y=1]
\drawface;
\node[rotate=-90] {$\overline{1}$};
\end{scope}
}
\newcommand\toponebar{
\begin{scope}[canvas is yx plane at z=1]
\drawface;
\node[yscale=-1] {$\overline{1}$};
\end{scope}
}
\newcommand\frontonebar{
\begin{scope}[canvas is yz plane at x=1]
\drawface;
\node  {$\overline{1}$};
\end{scope}
}
\newcommand\rightone{
\begin{scope}[canvas is zx plane at y=1]
\drawface;
\node[rotate=-90] {$1$};
\end{scope}
}
\newcommand\topone{
\begin{scope}[canvas is yx plane at z=1]
\drawface;
\node[yscale=-1] {$1$};
\end{scope}
}
\newcommand\frontone{
\begin{scope}[canvas is yz plane at x=1]
\drawface;
\node  {$1$};
\end{scope}
}
\newcommand\righttwo{
\begin{scope}[canvas is zx plane at y=1]
\drawface;
\node[rotate=-90] {$2$};
\end{scope}
}
\newcommand\toptwo{
\begin{scope}[canvas is yx plane at z=1]
\drawface;
\node[yscale=-1] {$2$};
\end{scope}
}
\newcommand\fronttwo{
\begin{scope}[canvas is yz plane at x=1]
\drawface;
\node  {$2$};
\end{scope}
}
\newcommand\rightthree{
\begin{scope}[canvas is zx plane at y=1]
\drawface;
\node[rotate=-90] {$3$};
\end{scope}
}
\newcommand\topthree{
\begin{scope}[canvas is yx plane at z=1]
\drawface;
\node[yscale=-1] {$3$};
\end{scope}
}
\newcommand\frontthree{
\begin{scope}[canvas is yz plane at x=1]
\drawface;
\node  {$3$};
\end{scope}
}
\newcommand\onebarcube[2]{
\begin{scope}[xshift=#1, yshift=#2]
\colorlet{heapcolor}{IBMpurple}
\faceone \facetwo \facethree \rightonebar \toponebar \frontonebar
\end{scope}
}
\newcommand\onecube[2]{
\begin{scope}[xshift=#1, yshift=#2]
\colorlet{heapcolor}{IBMblue}
\faceone \facetwo \facethree \rightone \topone \frontone
\end{scope}
}
\newcommand\twocube[2]{
\begin{scope}[xshift=#1, yshift=#2]
\colorlet{heapcolor}{IBMgold}
\faceone \facetwo \facethree \righttwo \toptwo \fronttwo
\end{scope}
}
\newcommand\threecube[2]{
\begin{scope}[xshift=#1, yshift=#2]
\colorlet{heapcolor}{IBMmagenta}
\faceone \facetwo \facethree \rightthree \topthree \frontthree
\end{scope}
}
\onebarcube{-14.25}{-42.75};
\twocube{0}{0};
\onecube{-42.75}{42.75};
\threecube{28.5}{57};
\twocube{0}{114};
\onebarcube{-14.25}{185.25};
\end{tikzpicture}
}
} \\[12pt]

\figurecell{Cayley diagram for $Q_8$}{%
% Copied from ProductsQuotients.tex, Figure 6.2(a).
{
\colorlet{cb-orange}{IBMmagenta}
\colorlet{cb-purple}{IBMblue}
\colorlet{cb-green}{IBMgold}
\tikzstyle{vert} = [quotient vertex]
\tikzstyle{p} = [draw, very thick, cb-purple,-stealth]
\tikzstyle{o} = [draw, very thick, cb-orange,-stealth]
\tikzstyle{g2} = [draw, very thick, cb-green, stealth-stealth]
\begin{tikzpicture}[scale=1.5,auto]
\node (1) at (135:2) [vert] {};
\node (i) at (45:2) [vert] {};
\node (k) at (-45:2) [vert] {};
\node (j) at (-135:2) [vert] {};
\node (-1) at (135:1) [vert] {};
\node (-i) at (45:1) [vert] {};
\node (-k) at (-45:1) [vert] {};
\node (-j) at (-135:1) [vert] {};
\path[p,figure edge] (1) to (i);
\path[p,figure edge] (i) to (-1);
\path[p,figure edge] (-1) to (-i);
\path[p,figure edge] (-i) to (1);
\path[p,figure edge] (-j) to (-k);
\path[p,figure edge] (-k) to (j);
\path[p,figure edge] (j) to (k);
\path[p,figure edge] (k) to (-j);
\path[o,figure edge] (-k) to (-i);
\path[o,figure edge] (-i) to (k);
\path[o,figure edge] (k) to (i);
\path[o,figure edge] (i) to (-k);
\path[o,figure edge] (1) to (j);
\path[o,figure edge] (j) to (-1);
\path[o,figure edge] (-1) to (-j);
\path[o,figure edge] (-j) to (1);
\path[g2,figure edge] (1) to (-1);
\path[g2,figure edge] (j) to (-j);
\path[g2,figure edge] (i) to (-i);
\path[g2,figure edge] (k) to (-k);
\end{tikzpicture}
}
} &
\figurecell{Structure diagram for a $2\times\text{odd}$ grid}{%
% Figure 7.3, left: Grid_2_3=GEN.pdf.
% Layout and shading reproduced from the published diagram.
{
\tikzset{
structure up/.style={draw,isosceles triangle,isosceles triangle apex angle=100,
shape border rotate=90,minimum width=17mm,minimum height=7mm,
inner sep=0pt,font=\small},
structure down/.style={structure up,shape border rotate=270},
structure arrow/.style={draw,thick,-stealth}
}
\begin{tikzpicture}[x=1cm,y=1cm]
\node[structure up] (T0) at (0.7,4.8) {$0,1$};
\node[structure down,colored fill,fill=IBMorange] (T1) at (1.4,3.6) {$0,2$};
\node[structure up,colored fill,fill=IBMorange] (T3) at (0,2.4) {$3,4$};
\node[structure up] (T7) at (2.8,2.4) {$0,2$};
\node[structure down,colored fill,fill=IBMorange] (T8) at (-1.4,1.2) {$2,1$};
\node[structure up,colored fill,fill=IBMorange] (T10) at (1.4,1.2) {$1,2$};
\node[structure up] (T12) at (0,0) {$0,0$};
\path[structure arrow,figure edge] (T0) edge (T1) edge[bend right=8] (T3);
\path[structure arrow,figure edge] (T1) edge (T3) edge (T7) edge (T10);
\path[structure arrow,figure edge] (T3) edge (T8) edge (T10) edge (T12);
\path[structure arrow,figure edge] (T7) edge (T10);
\path[structure arrow,figure edge] (T8) edge (T12);
\path[structure arrow,figure edge] (T10) edge (T12);
\end{tikzpicture}
}
} \\[12pt]

\figurecell{Representative quotient for $\mathrm{GEN}(W_5)$}{%
% Copied from GENandDNG-version3.tex, Figure 3.2, right.
{
\definecolor{grey}{rgb}{.7,.7,.7}
\colorlet{purple}{IBMpurple}
\tikzstyle{small vert} = [quotient vertex]
\tikzstyle{b} = [draw,very thick,black,-]
\tikzstyle{a} = [draw,black,-stealth]
\begin{tikzpicture}[rotate=-90,xscale=.85]
\node[inner sep=0pt] (middle1) at (0,0) {
\begin{tikzpicture}[scale=1,auto]
\node (1) at (0,0) [small vert] {};
\node (2) at (1,0) [small vert] {};
\node (3) at (1,1) [small vert] {};
\node (4) at (0,1) [small vert] {};
\node (5) at (.5,.5) [small vert] {};
\path[b,figure edge] (1) to (2) to (3) to (4) to (1);
\path[b,figure edge] (1) to (5);
\path[b,figure edge] (2) to (5);
\path[b,figure edge] (3) to (5);
\path[b,figure edge] (4) to (5);
\begin{pgfonlayer}{background}
\fill[white,opacity=0.3] \convexpath{4,3,2,1}{7pt};
\end{pgfonlayer}
\end{tikzpicture}
};
\node[inner sep=0pt] (middle2) at (2.25,0) {
\begin{tikzpicture}[scale=1,auto]
\filldraw[purple,opacity=0.3] (0,1) circle (7pt);
\node (1) at (0,0) [small vert] {};
\node (2) at (1,0) [small vert] {};
\node (3) at (1,1) [small vert] {};
\node (4) at (0,1) [small vert,fill=purple] {};
\node (5) at (.5,.5) [small vert] {};
\path[b,figure edge] (1) to (2) to (3) to (4) to (1);
\path[b,figure edge] (1) to (5);
\path[b,figure edge] (2) to (5);
\path[b,figure edge] (3) to (5);
\path[b,figure edge] (4) to (5);
% \begin{pgfonlayer}{background}
% \fill[white,opacity=0.3] \convexpath{4,4}{7pt};
% \end{pgfonlayer}
\begin{pgfonlayer}{background}
\fill[white,opacity=0.3] \convexpath{4,3,2,1}{7pt};
\end{pgfonlayer}
\end{tikzpicture}
};
\node[inner sep=0pt] (right2) at (2.25,2.25) {
\begin{tikzpicture}[scale=1,auto]
\filldraw[purple,opacity=0.3] (.5,.5) circle (7pt);
\node (1) at (0,0) [small vert] {};
\node (2) at (1,0) [small vert] {};
\node (3) at (1,1) [small vert] {};
\node (4) at (0,1) [small vert] {};
\node (5) at (.5,.5) [small vert,fill=purple] {};
\path[b,figure edge] (1) to (2) to (3) to (4) to (1);
\path[b,figure edge] (1) to (5);
\path[b,figure edge] (2) to (5);
\path[b,figure edge] (3) to (5);
\path[b,figure edge] (4) to (5);
% \begin{pgfonlayer}{background}
% \fill[white,opacity=0.3] \convexpath{4,4}{7pt};
% \end{pgfonlayer}
\begin{pgfonlayer}{background}
\fill[white,opacity=0.3] \convexpath{4,3,2,1}{7pt};
\end{pgfonlayer}
\end{tikzpicture}
};
\node[inner sep=0pt] (left3) at (4.5,-2.25) {
\begin{tikzpicture}[scale=1,auto]
\node (1) at (0,0) [small vert] {};
\node (2) at (1,0) [small vert,fill=purple] {};
\node (3) at (1,1) [small vert] {};
\node (4) at (0,1) [small vert,fill=purple] {};
\node (5) at (.5,.5) [small vert] {};
\path[b,figure edge] (1) to (2) to (3) to (4) to (1);
\path[b,figure edge] (1) to (5);
\path[b,figure edge] (2) to (5);
\path[b,figure edge] (3) to (5);
\path[b,figure edge] (4) to (5);
\begin{pgfonlayer}{background}
\fill[purple,opacity=0.3] \convexpath{1,4,3,2}{7pt};
\end{pgfonlayer}
\end{tikzpicture}
};
\node[inner sep=0pt] (middle3) at (4.5,0) {
\begin{tikzpicture}[scale=1,auto]
\node (1) at (0,0) [small vert,fill=purple] {};
\node (2) at (1,0) [small vert] {};
\node (3) at (1,1) [small vert] {};
\node (4) at (0,1) [small vert,fill=purple] {};
\node (5) at (.5,.5) [small vert] {};
\path[b,figure edge] (1) to (2) to (3) to (4) to (1);
\path[b,figure edge] (1) to (5);
\path[b,figure edge] (2) to (5);
\path[b,figure edge] (3) to (5);
\path[b,figure edge] (4) to (5);
\begin{pgfonlayer}{background}
\fill[purple,opacity=0.3] \convexpath{1,4}{7pt};
\end{pgfonlayer}
\begin{pgfonlayer}{background}
\fill[white,opacity=0.3] \convexpath{4,3,2,1}{7pt};
\end{pgfonlayer}
\end{tikzpicture}
};
\node[inner sep=0pt] (right3) at (4.5,2.25) {
\begin{tikzpicture}[scale=1,auto]
\node (1) at (0,0) [small vert] {};
\node (2) at (1,0) [small vert] {};
\node (3) at (1,1) [small vert] {};
\node (4) at (0,1) [small vert,fill=purple] {};
\node (5) at (.5,.5) [small vert,fill=purple] {};
\path[b,figure edge] (1) to (2) to (3) to (4) to (1);
\path[b,figure edge] (1) to (5);
\path[b,figure edge] (2) to (5);
\path[b,figure edge] (3) to (5);
\path[b,figure edge] (4) to (5);
\begin{pgfonlayer}{background}
\fill[purple,opacity=0.3] \convexpath{4,5}{7pt};
\end{pgfonlayer}
\begin{pgfonlayer}{background}
\fill[white,opacity=0.3] \convexpath{4,3,2,1}{7pt};
\end{pgfonlayer}
\end{tikzpicture}
};
\node[inner sep=0pt] (middle4) at (6.75,0) {
\begin{tikzpicture}[scale=1,auto]
\node (1) at (0,0) [small vert,fill=purple] {};
\node (2) at (1,0) [small vert] {};
\node (3) at (1,1) [small vert] {};
\node (4) at (0,1) [small vert,fill=purple] {};
\node (5) at (.5,.5) [small vert,fill=purple] {};
\path[b,figure edge] (1) to (2) to (3) to (4) to (1);
\path[b,figure edge] (1) to (5);
\path[b,figure edge] (2) to (5);
\path[b,figure edge] (3) to (5);
\path[b,figure edge] (4) to (5);
\begin{pgfonlayer}{background}
\fill[purple,opacity=0.3] \convexpath{1,4,5}{7pt};
\end{pgfonlayer}
\begin{pgfonlayer}{background}
\fill[white,opacity=0.3] \convexpath{4,3,2,1}{7pt};
\end{pgfonlayer}
\end{tikzpicture}
};
\node[inner sep=0pt] (left4) at (6.75,-2.25) {
\begin{tikzpicture}[scale=1,auto]
\node (1) at (0,0) [small vert,fill=purple] {};
\node (2) at (1,0) [small vert,fill=purple] {};
\node (3) at (1,1) [small vert] {};
\node (4) at (0,1) [small vert,fill=purple] {};
\node (5) at (.5,.5) [small vert] {};
\path[b,figure edge] (1) to (2) to (3) to (4) to (1);
\path[b,figure edge] (1) to (5);
\path[b,figure edge] (2) to (5);
\path[b,figure edge] (3) to (5);
\path[b,figure edge] (4) to (5);
\begin{pgfonlayer}{background}
\fill[purple,opacity=0.3] \convexpath{1,4,3,2}{7pt};
\end{pgfonlayer}
\end{tikzpicture}
};
\node[inner sep=0pt] (right4) at (6.75,2.25) {
\begin{tikzpicture}[scale=1,auto]
\node (1) at (0,0) [small vert] {};
\node (2) at (1,0) [small vert,fill=purple] {};
\node (3) at (1,1) [small vert] {};
\node (4) at (0,1) [small vert,fill=purple] {};
\node (5) at (.5,.5) [small vert,fill=purple] {};
\path[b,figure edge] (1) to (2) to (3) to (4) to (1);
\path[b,figure edge] (1) to (5);
\path[b,figure edge] (2) to (5);
\path[b,figure edge] (3) to (5);
\path[b,figure edge] (4) to (5);
\begin{pgfonlayer}{background}
\fill[purple,opacity=0.3] \convexpath{1,4,3,2}{7pt};
\end{pgfonlayer}
\end{tikzpicture}
};
\node[inner sep=0pt] (middle5) at (9,0) {
\begin{tikzpicture}[scale=1,auto]
\node (1) at (0,0) [small vert,fill=purple] {};
\node (2) at (1,0) [small vert] {};
\node (3) at (1,1) [small vert,fill=purple] {};
\node (4) at (0,1) [small vert,fill=purple] {};
\node (5) at (.5,.5) [small vert,fill=purple] {};
\path[b,figure edge] (1) to (2) to (3) to (4) to (1);
\path[b,figure edge] (1) to (5);
\path[b,figure edge] (2) to (5);
\path[b,figure edge] (3) to (5);
\path[b,figure edge] (4) to (5);
\begin{pgfonlayer}{background}
\fill[purple,opacity=0.3] \convexpath{1,4,3,2}{7pt};
\end{pgfonlayer}
\end{tikzpicture}
};
\path[a,figure edge] (middle1) to (middle2);
\path[a,figure edge] (middle1) to (right2);
\path[a,figure edge] (middle2) to (left3);
\path[a,figure edge] (middle2) to (middle3);
\path[a,figure edge] (middle2) to (right3);
\path[a,figure edge] (right2) to (right3);
\path[a,figure edge] (middle3) to (left4);
\path[a,figure edge] (middle3) to (middle4);
\path[a,figure edge] (middle4) to (middle5);
\path[a,figure edge] (right3) to (right4);
\path[a,figure edge] (right3) to (middle4);
\end{tikzpicture}
}
} &
\figurecell{Temperley--Lieb diagram factorization}{%
% ANAS-Salmon.tex, slide 8, final overlay, right subfigure.
{
\colorlet{butterfly}{IBMblue}
\colorlet{rose}{IBMmagenta}
\colorlet{mediterranean}{IBMpurple}
\colorlet{darkorchid}{IBMorange}
\begin{tikzpicture}[scale=.45]
%d_1
\begin{scope}
\clip{(1,0) rectangle (1.855,3)};
\fill[darkorchid](1,3)--(1.856,1.71)--(1.856,0)--(1,0);
\fill[white] (1,3) to [out=-80,in=90]  (3,0);
\end{scope}
\begin{scope}
\clip{(1.85,0) rectangle (2,3)};
\fill[darkorchid](2,1.55)--(1.85,1.71)--(1.85,0)--(2,0);
\fill[darkorchid] (1,3) to [out=-80,in=90]  (3,0);
\end{scope}
\begin{scope}
\clip{(1,0) rectangle (2,3)};
\fill[white](1,0) to [bend left=60] (2,0);
\end{scope}
%d_2
\begin{scope}
\clip{(2,0) rectangle (3,3)};
\fill[mediterranean](2,1.57)--(3,0)--(3,1.4)--(2,3);
\fill[white] (1,3) to [out=-80,in=90]  (3,0);
\fill[white](2,3) to [out=-80,in=90]  (4,0);
\fill[mediterranean](2,1.7)--(3,1.3)--(3,1.58)--(2,2.45);
\end{scope}
%d_3
\begin{scope}
\clip{(3,0) rectangle (4,3)};
\fill[rose](3,1.57)--(4,0)--(4,3)--(3,3);
\fill[white] (3,3) to [bend right=60] (6, 3);
\fill[white](2,3) to [out=-80,in=90]  (4,0);
\end{scope}
%d_4
\begin{scope}
\clip{(4,0) rectangle (5,3)};
\fill[butterfly](3,3) to [bend right=60] (6, 3);
\fill[white](4,3) to [bend right=60] (5,3);
\end{scope}
%d_5
\begin{scope}
\clip{(5,0) rectangle (6,3)};
\fill[rose](5,0)--(6,.9)--(6,3)--(5,3);
\fill[white] (3,3) to [bend right=60] (6, 3);
\fill[white](9,3) to [out=-110,in=60]  (5,0);
\end{scope}
%d_6
\begin{scope}
\clip{(6,0) rectangle (6.74,3)};
\fill[mediterranean](9,3) to [out=-110,in=60]  (5,0);
\fill[mediterranean](6,0)--(6.76,0) -- (6.76,1.315)--(6,.8);
%\fill[white] (6,0) to [bend left=60] (9,0);
\end{scope}
\begin{scope}
\clip{(6.74,0) rectangle (7,3)};
\fill[mediterranean](6,0)--(7,0) -- (7,1.48)--(6,.8);
\fill[white](9,3) to [out=-110,in=60]  (5,0);
%\fill[white] (6,0) to [bend left=60] (9,0);
\end{scope}
\begin{scope}
\clip{(6,0) rectangle (7,3)};
\fill[white] (6,0) to [bend left=60] (9,0);
\end{scope}
%d_7
\begin{scope}
\clip{(7,0) rectangle (8,3)};
\fill[butterfly](7,1.5)--(8,1.9) -- (8,3) -- (7,3);
\fill[butterfly](9,3) to [out=-110,in=60]  (5,0);
\fill[white] (7,3) to [bend right=60] (8,3);
\fill[darkorchid](6,0) to [bend left=60] (9,0);
\fill[white] (7,0) to [bend left=60] (8,0);
\end{scope}
%d_8
\begin{scope}
\clip{(8,0) rectangle (9,3)};
\fill[rose](8,0)--(9,0)--(9,3)--(8,1.9);
\fill[white] (6,0) to [bend left=60] (9,0);
\fill[white](9,3) to [out=-110,in=60]  (5,0);
\end{scope}
\draw[gray, thick,figure edge] (0,0) rectangle (10,3);
\foreach \x in {1,2,3,4,5,6,7,8,9} \filldraw (\x,0) circle (1pt);
\foreach \x in {1,2,3,4,5,6,7,8,9} \filldraw (\x,3) circle (1pt);
\draw[figure edge](1,0) to [bend left=60] (2,0);
\draw[figure edge](3,3) to [bend right=60] (6, 3);
\draw[figure edge](4,3) to [bend right=60] (5,3);
\draw[figure edge](6,0) to [bend left=60] (9,0);
\draw[figure edge](7,0) to [bend left=60] (8,0);
\draw[figure edge](1,3) to [out=-80,in=90]  (3,0);
\draw[figure edge](2,3) to [out=-80,in=90]  (4,0);
\draw[figure edge](9,3) to [out=-110,in=60]  (5,0);
\draw[figure edge](7,3) to [bend right=60] (8,3);
%
%vertical lines
\draw[dashed,figure edge] (1,0) to (1,3);
\draw[dashed,figure edge] (2,0) to (2,3);
\draw[dashed,figure edge] (3,0) to (3,3);
\draw[dashed,figure edge] (4,0) to (4,3);
\draw[dashed,figure edge] (5,0) to (5,3);
\draw[dashed,figure edge] (6,0) to (6,3);
\draw[dashed,figure edge] (7,0) to (7,3);
\draw[dashed,figure edge] (8,0) to (8,3);
\draw[dashed,figure edge] (9,0) to (9,3);
\end{tikzpicture}
}
} \\[12pt]

\figurecell{Nim gamegraph for piles of sizes 3 and 2}{%
% Copied from main.tex, Figure 3, left subfigure.
{
\tikzstyle{rect vert} = [rectangle,rounded corners,draw,fill=grey!40,inner sep=2.7pt,minimum size=4mm]
\tikzstyle{d} = [draw,thick,black,-stealth]
\definecolor{grey}{rgb}{.7, .7, .7}
\colorlet{purple}{IBMpurple}
\colorlet{orange}{IBMgold}
\colorlet{turq}{IBMgold}
\colorlet{springgreen}{IBMblue!70}
\colorlet{rose}{IBMmagenta}
\colorlet{darkorchid}{IBMorange}
\colorlet{mediterranean}{IBMpurple}
\colorlet{rred}{IBMmagenta}
\begin{tikzpicture}[scale=1.2,auto]
\node (1) at (0,5) [rect vert,fill=purple!50] {\scriptsize $(3,2)$};
\node (2) at (-1,4) [rect vert,fill=orange!50] {\scriptsize $(2,2)$};
\node (3) at (1,4) [rect vert,fill=turq!50] {\scriptsize $(3,1)$};
\node (4) at (-2,3) [rect vert,fill=springgreen!50] {\scriptsize $(1,2)$};
\node (5) at (0,3) [rect vert,fill=springgreen!50] {\scriptsize $(2,1)$};
\node (6) at (2,3) [rect vert,fill=rose!50] {\scriptsize $(3,0)$};
\node (7) at (-2,2) [rect vert,fill=darkorchid!50] {\scriptsize $(1,1)$};
\node (8) at (0,2) [rect vert] {\scriptsize $(0,2)$};
\node (9) at (2,2) [rect vert] {\scriptsize $(2,0)$};
\node (10) at (-1,1) [rect vert,fill=mediterranean!50] {\scriptsize $(1,0)$};
\node (11) at (1,1) [rect vert,fill=mediterranean!50] {\scriptsize$(0,1)$};
\node (12) at (0,0) [rect vert,fill=rred!50] {\scriptsize $(0,0)$};
%options 1
\path [d,out=196,in=80] (1) to (4);
\path [d] (1) to (2);
\path [d,out=-60,in=45] (1) to (8);
\path [d] (1) to (3);
\path [d,out=-15,in=110] (1) to (6);
%options 2
\path [d] (2) to (4);
\path [d,out=-90] (2) to (8);
\path [d] (2) to (5);
\path [d,out=-20] (2) to (9);
%options 3
\path [d,out=200,in=45] (3) to (7);
\path [d] (3) to (5);
\path [d] (3) to (6);
\path [d] (3) to (11);
%options 4
\path [d,out=-45,in=90] (4) to (10);
\path [d] (4) to (7);
\path [d] (4) to (8);
%options 5
\path [d] (5) to (7);
\path [d,out=-45,in=90] (5) to (11);
\path [d] (5) to (9);
%options 6
\path [d,out=-130,in=30] (6) to (10);
\path [d] (6) to (9);
\path [d,out=-40,in=0] (6) to (12);
%options 7
\path [d] (7) to (10);
\path [d,out=-20,in=165] (7) to (11);
%options 8
\path [d] (8) to (12);
\path [d] (8) to (11);
%options 9
\path [d,out=210,in=10] (9) to (10);
\path [d,out=-95,in=20] (9) to (12);
%options 10
\path [d] (10) to (12);
%options 11
\path [d] (11) to (12);
\end{tikzpicture}
}
} &
\figurecell{Gamegraph on subsets of $\{1,2,3,4\}$}{%
\begin{tikzpicture}[xscale=2.1,yscale=1.1,inner sep=.3mm]
%\draw[line width=7mm,yellow!55,rounded corners, line cap=round] (0,3) -- (.4,2.5) -- (.4,1.5) -- (0,1);
%\draw[line width=7mm,yellow!55,rounded corners, line cap=round] (-.5,2) -- (0,2);
%\draw[line width=7mm,yellow!55,rounded corners, line cap=round] (0-1.5,3-1) -- (.4-1.5,2.5-1) -- (.4-1.5,1.5-1) -- (0-1.5,1-1);
%\draw[line width=7mm,yellow!55,rounded corners, line cap=round] (-.5-1.5,2-1) -- (0-1.5,2-1);
%\draw[line width=7mm,yellow!55,rounded corners, line cap=round] (0+1.5,3-1) -- (.4+1.5,2.5-1) -- (.4+1.5,1.5-1) -- (0+1.5,1-1);
%\draw[line width=7mm,yellow!55,rounded corners, line cap=round] (-.5+1.5,2-1) -- (0+1.5,2-1);
\node[rect vert,colored fill,fill=green] (0) at (0,3) {$\scriptstyle\emptyset$};
\node[rect vert,colored fill,fill=orange] (1) at (-1.5,1.5) {${\scriptstyle\{1\}}$};
\node[rect vert,fill=green!40] (2) at (-.3,2) {${\scriptstyle\{2\}}$};
\node[rect vert,fill=green!40] (3) at (.3,2) {${\scriptstyle\{3\}}$};
\node[rect vert,colored fill,fill=red] (4) at (1.5,1.5) {${\scriptstyle\{4\}}$};
\node[rect vert,fill=orange!40] (12) at (-1.8,0.5) {${\scriptstyle\{1,2\}}$};
\node[rect vert,fill=orange!40] (13) at (-1.2,0.5) {${\scriptstyle\{1,3\}}$};
\node[rect vert,colored fill,fill=green] (23) at (0,1) {${\scriptstyle\{2,3\}}$};
\node[rect vert,fill=red!40] (24) at (1.2,.5) {${\scriptstyle\{2,4\}}$};
\node[rect vert,fill=red!40] (34) at (1.8,.5) {${\scriptstyle\{3,4\}}$};
\node[rect vert,colored fill,fill=orange](123) at (-1.5,-.5) {${\scriptstyle\{1,2,3\}}$};
\node[rect vert,colored fill,fill=red] (234) at (1.5,-.5) {${\scriptstyle\{2,3,4\}}$};
\node[rect vert,colored fill,fill=blue] (14) at (0,0) {${\scriptstyle\{1,4\}}$};
\node[rect vert,fill=blue!40] (124) at (-.4,-1) {${\scriptstyle\{1,2,4\}}$};
\node[rect vert,fill=blue!40] (134) at (.4,-1) {${\scriptstyle\{1,3,4\}}$};
\node[rect vert,colored fill,fill=blue] (1234) at (0,-2) {${\scriptstyle\{1,2,3,4\}}$};
\draw[thick,densely dotted,d,figure edge] (0)--(3);
\draw[thick,densely dotted,d,figure edge] (0)--(2);
\draw[d,figure edge] (0)--(1);
\draw[d,figure edge] (0)--(4);
\draw[thick,densely dotted,d,figure edge] (1)--(12);
\draw[thick,densely dotted,d,figure edge] (1)--(13);
\draw[d,figure edge] (2)--(23);
\draw[d,figure edge] (2)--(12);
\draw[d,figure edge] (2)--(24);
\draw[d,figure edge] (3)--(13);
\draw[d,figure edge] (3)--(23);
\draw[d,figure edge] (3)--(34);
\draw[thick,densely dotted,d,figure edge] (4)--(24);
\draw[thick,densely dotted,d,figure edge] (4)--(34);
\draw[d,figure edge] (12)--(123);
\draw[d,figure edge] (13)--(123);
\draw[d,figure edge] (23)--(123);
\draw[d,figure edge] (23)--(234);
\draw[d,figure edge] (24)--(234);
\draw[d,figure edge] (34)--(234);
\draw[d,figure edge] (1)--(14);
\draw[d,figure edge] (4)--(14);
\draw[d,figure edge] (12)--(124);
\draw[d,figure edge] (24)--(124);
\draw[d,figure edge] (13)--(134);
\draw[d,figure edge] (34)--(134);
\draw[d,figure edge] (123)--(1234);
\draw[d,figure edge] (234)--(1234);
\end{tikzpicture}
} \\[12pt]

\figurecell{L-tromino tiling}{%
% IntroInduction.tex, Figure 4.1.
{
\definecolor{grey2}{RGB}{204,204,204}
\definecolor{grey3}{RGB}{150,150,150}
\definecolor{grey5}{RGB}{37,37,37}
\begin{tikzpicture} [scale=1.5]
\draw [tromino border,colored fill,fill=IBMblue] (0,0) rectangle (2,1);
\draw [tromino border,colored fill,fill=IBMblue] (0,1) rectangle (1,2);
\draw [tromino border,colored fill,fill=IBMmagenta] (3,0) rectangle (4,2);
\draw [tromino border,colored fill,fill=IBMmagenta] (2,1) rectangle (3,2);
\draw [tromino border,colored fill,fill=IBMorange] (0,2) rectangle (1,4);
\draw [tromino border,colored fill,fill=IBMorange] (1,3) rectangle (2,4);
\draw [tromino border,colored fill,fill=IBMgold] (1,1) rectangle (2,3);
\draw [tromino border,colored fill,fill=IBMgold] (2,2) rectangle (3,3);
\draw [tromino border,colored fill,fill=IBMblue] (2,3) rectangle (4,4);
\draw [tromino border,colored fill,fill=IBMblue] (3,2) rectangle (4,3);
\draw[tromino border] (0,0) grid (4,4);
\end{tikzpicture}
}
} &
\figurecell{Function diagram with fibers}{%
% InjectiveSurjectiveFunctions.tex, Figure 8.2, left subfigure.
{
\definecolor{grey}{rgb}{.6,.6,.6}
\tikzstyle{vert} = [circle,draw,fill=grey,inner sep=0pt,minimum size=5mm]
\begin{tikzpicture}[->,>=stealth,shorten >=1pt,auto,node distance=2cm,semithick,scale=1.2]
\node[vert] (A) at (0,6) {\scriptsize $a$};
\node[vert] (B) at (0,5) {\scriptsize $b$};
\node[vert] (C) at (0,4) {\scriptsize $c$};
\node[vert] (D) at (0,3) {\scriptsize $d$};
\node[vert] (E) at (0,2) {\scriptsize $e$};
\node[vert] (F) at (0,1) {\scriptsize $f$};
\node[vert] (1) at (2.5,5.5) {\scriptsize $1$};
\node[vert] (2) at (2.5,4.5) {\scriptsize $2$};
\node[vert] (3) at (2.5,3.5) {\scriptsize $3$};
\node[vert] (4) at (2.5,2.5) {\scriptsize $4$};
\node[vert] (5) at (2.5,1.5) {\scriptsize $5$};
\path (A) edge node {} (1)
(B) edge node {} (1)
(C) edge node {} (2)
(D) edge node {} (4)
(E) edge node {} (4)
(F) edge node {} (4);
\node[draw,rounded corners,fit=(A) (B) (C) (D) (E) (F),minimum width=1.3cm,minimum height=7.25cm,gray] {};
\node[draw,rounded corners,fit=(1) (2) (3) (4) (5),minimum width=1.3cm,gray] {};
\node[draw,rounded corners,fit=(A) (B),minimum width=.9cm,gray] {};
\node[draw,rounded corners,fit=(C),minimum width=.9cm,gray] {};
\node[draw,rounded corners,fit=(D) (E) (F),minimum width=.9cm,gray] {};
\node (X) at (0,6.75) {$X$};
\node (Y) at (2.5,6.1) {$Y$};
\path (X) edge [bend left] node {$\varphi$} (Y);
\end{tikzpicture}
}
} \\[12pt]

\figurecell{Matsumoto graph}{%
% ---------------- BOTTOM: MAIN DIAGRAM ----------------
\begin{tikzpicture}[
every circle node/.style={graph vertex},
scale=1.1
]
% -------- LEFT GRAPH --------
\begin{scope}[local bounding box=leftbox]
\node[circle] (a) at (1,0) {};
\node[circle] (b) at (2,0.5) {};
\node[circle] (c) at (3,0.5) {};
\node[circle] (d) at (3,1) {};
\node[circle] (e) at (4,1) {};
\node[circle] (f) at (5,1.5) {};
\node[circle] (g) at (0,1) {};
\node[circle] (h) at (1,1.5) {};
\node[circle] (i) at (2,2) {};
\node[circle] (j) at (0,2) {};
\node[circle] (k) at (1,2.5) {};
\node[circle] (l) at (2,3) {};
\node[circle] (m) at (4,1.5) {};
\node[circle] (n) at (5,2) {};
\node[circle] (o) at (6,2.5) {};
\node[circle] (p) at (4,2.5) {};
\node[circle] (q) at (5,3) {};
\node[circle] (r) at (6,3.5) {};
\node[circle] (s) at (2,3.5) {};
\node[circle] (t) at (3,3.5) {};
\node[circle] (u) at (3,4) {};
\node[circle] (v) at (4,4) {};
% Background Layer
\begin{pgfonlayer}{background}
\draw[line width=8pt, blue, line cap=round, line join=round, opacity=0.25,figure edge]
(a.center) -- (b.center) -- (d.center) -- (e.center) -- (f.center) -- (c.center) -- (b.center);
\end{pgfonlayer}
\draw[orange, very thick,figure edge] (l) -- (u);
% Blue Edges
\draw[blue, very thick,figure edge] (a) -- (b);
\draw[blue, very thick,figure edge] (b) -- (c);
\draw[blue, very thick,figure edge] (b) -- (d);
\draw[blue, very thick,figure edge] (c) -- (e);
\draw[blue, very thick,figure edge] (d) -- (e);
\draw[blue, very thick,figure edge] (e) -- (f);
\draw[blue, very thick,figure edge] (g) -- (h);
\draw[blue, very thick,figure edge] (h) -- (i);
\draw[blue, very thick,figure edge] (j) -- (k);
\draw[blue, very thick,figure edge] (k) -- (l);
\draw[blue, very thick,figure edge] (m) -- (n);
\draw[blue, very thick,figure edge] (n) -- (o);
\draw[blue, very thick,figure edge] (p) -- (q);
\draw[blue, very thick,figure edge] (q) -- (r);
\draw[blue, very thick,figure edge] (s) -- (u);
\draw[blue, very thick,figure edge] (v) -- (u);
\draw[blue, very thick,figure edge] (v) -- (t);
\draw[blue, very thick,figure edge] (s) -- (t);
% Orange Edges
\draw[orange, very thick,figure edge] (a) -- (g);
\draw[orange, very thick,figure edge] (g) -- (j);
\draw[orange, very thick,figure edge] (b) -- (h);
\draw[orange, very thick,figure edge] (h) -- (k);
\draw[orange, very thick,figure edge] (k) -- (s);
\draw[orange, very thick,figure edge] (d) -- (i);
\draw[orange, very thick,figure edge] (i) -- (l);
\draw[orange, very thick,figure edge] (c) -- (m);
\draw[orange, very thick,figure edge] (m) -- (p);
\draw[orange, very thick,figure edge] (p) -- (t);
\draw[orange, very thick,figure edge] (e) -- (n);
\draw[orange, very thick,figure edge] (n) -- (q);
\draw[orange, very thick,figure edge] (q) -- (v);
\draw[orange, very thick,figure edge] (f) -- (o);
\draw[orange, very thick,figure edge] (o) -- (r);
\end{scope}
\end{tikzpicture}
} &
 \\[12pt]
\end{longtable}
\end{document}
